% Generated draft: requires mathematical and source review.
\documentclass[11pt]{article}
\usepackage[margin=25mm]{geometry}
\usepackage{fontspec}
\setmainfont{Latin Modern Roman}
\usepackage{amsmath,amssymb,mathtools,unicode-math}
\setmathfont{Latin Modern Math}
\usepackage{enumitem,needspace,xurl,hyperref}
\setlength{\parindent}{0pt}
\setlength{\parskip}{5pt}
\setlength{\emergencystretch}{3em}
\newcommand{\N}{\mathbb{N}}
\newcommand{\Z}{\mathbb{Z}}
\newcommand{\Q}{\mathbb{Q}}
\newcommand{\R}{\mathbb{R}}
\newcommand{\C}{\mathbb{C}}
\newcommand{\F}{\mathbb{F}}
\newcommand{\A}{\mathbb{A}}
\newcommand{\PP}{\mathbb{P}}
\newcommand{\E}{\mathbb{E}}
\newcommand{\eps}{\varepsilon}
\newcommand{\dd}{\,\mathrm{d}}
\begin{document}
\section*{UnsolvedMath ID: 700391}
\textbf{Source problem code:} catalogue entry
\subsection{Definitions and mathematical statement}
Let \(p\) be a prime number. A \(p\)-group here means a group \(G\) such that every element has finite order equal to a power of \(p\):
\[
\forall g\in G\;\exists n\in\Z_{\geq 0}\quad \operatorname{ord}(g)=p^n.
\]
Interpret \emph{binary finite} as meaning that every subgroup generated by two elements is finite:
\[
\forall x,y\in G\quad |\langle x,y\rangle|<\infty.
\]
The supplied record does not explicitly define this terminology; the formulation below uses this interpretation and requires confirmation against the cited source.

\textbf{Question (N. S. Chernikov, 1978; Kourovka Notebook, Problem 6.48).} For every prime \(p\) and every infinite binary finite \(p\)-group \(G\), must there exist a normal subgroup \(N\) satisfying
\[
\{1_G\}\subsetneq N\subsetneq G?
\]
Equivalently, must every such group be non-simple?
\subsection{Short English statement}
Chernikov asks whether, for every prime \(p\), an infinite \(p\)-group in which every two elements generate a finite subgroup must have a nontrivial proper normal subgroup. This formulation uses the stated interpretation of ``binary finite,'' which the supplied record leaves undefined.
\subsection{Sources}
\begin{itemize}
\item Local UnsolvedMath registry, permanent ID \(700391\), source code \texttt{catalogue entry}, title ``Kourovka Notebook Problem 6.48.'' The supplied record attributes the problem to N. S. Chernikov (1978) and labels it open; this status is supplied metadata, not an independently verified mathematical assessment.
\item \emph{The Kourovka Notebook (Unsolved Problems in Group Theory)}, original Problem 6.48, 6th issue (1978), as identified in the supplied record through the 21st issue edition (2026), \texttt{arXiv:1401.0300v48}. The collection metadata names E. I. Khukhro and V. D. Mazurov as editors. This reference is supplied attribution and has not been independently consulted.
\end{itemize}
\end{document}
